\section{Strutturazione Semantica dei Documenti tramite LLM}
\label{sec:strutturazione_semantica}

Una volta ultimata la fase di estrazione deterministica del testo grezzo, illustrata nella sezione precedente, i documenti normativi si presentano sotto forma di file testuali piatti in formato Markdown. Tali file, pur contenendo l'intero testo originale, sono del tutto privi di metadati strutturali e di una formattazione semantica chiara. 
Poiché l'architettura di sistema prevede l'ingestione di questi documenti all'interno di un database a grafo (Neo4j) e in un database vettoriale (Qdrant), risulta indispensabile trasformare il testo non strutturato in una gerarchia dati rigorosa. L'archiviazione del solo testo grezzo, infatti, risulterebbe del tutto inadeguata e priva di utilità pratica. Una strutturazione accurata rappresenta il prerequisito fondamentale per abilitare qualsiasi operazione avanzata di \textit{information retrieval}, garantendo che i dati possano essere interrogati in modo efficiente e contestuale.

\subsection{Il Modello e la Strategia di Generazione}

Per affrontare il compito cruciale della strutturazione sintattica, è stato impiegato il Large Language Model \texttt{Gemini 2.5-flash}. A livello di architettura software, l'integrazione del modello è stata realizzata mediante il modulo \texttt{JsonStrategy}, una classe scritta in Python che implementa il pattern comportamentale \textit{Strategy} interfacciandosi con le API di Google GenAI.

Al fine di garantire che l'output del modello sia strettamente confinato a una struttura dati ben definita, la strategia sfrutta la configurazione nativa delle API di Gemini, forzando il parametro \texttt{response\_mime\_type} al valore \texttt{"application/json"}. Questo vincolo formale assicura che il modello non restituisca testo discorsivo, introduzioni o blocchi di codice Markdown, ma esclusivamente un oggetto JSON puro, requisito fondamentale per la successiva automazione delle pipeline di inserimento a database.

\subsection{Mappatura dello Schema}

L'efficacia e l'accuratezza del processo di strutturazione si fondano su un \textit{System Prompt} ingegnerizzato ad hoc. Il modello viene calato nel ruolo di un esperto analista giuridico, con la direttiva assoluta di non riassumere o alterare in alcun modo il significato giuridico del testo originale. Il compito esclusivo del LLM è quello di sanare i difetti di impaginazione (come i ritorni a capo spuri derivanti dall'estrazione OCR) e mappare il contenuto all'interno del seguente schema logico:

\begin{itemize}
    \item \textbf{Metadati (\texttt{metadata}):} Il modello estrae e categorizza i dati identificativi dell'atto. Vengono isolati campi chiave come la data di stipula, il titolo completo e il nome del file originale.
    \item \textbf{Sezioni Introduttive e Conclusive:} Vengono separate dal corpo normativo tutte le sezioni di contorno. Tra queste figurano le premesse iniziali e le epigrafi (\texttt{preface}), gli indici (\texttt{preamble}), le dichiarazioni di chiusura con le relative firme (\texttt{conclusions}) e gli eventuali documenti annessi (\texttt{attachments}).
    \item \textbf{Corpo del Documento (\texttt{body}):} Rappresenta il nucleo dell'algoritmo di strutturazione. Il modello deve nidificare il testo seguendo fedelmente la gerarchia complessa e rigida tipica degli atti normativi italiani: \textbf{Titoli $>$ Capi $>$ Articoli $>$ Commi}.
\end{itemize}

Per garantire la massima coerenza dei dati all'interno del Knowledge Graph, sono state imposte al LLM regole tassative di parsing:

\begin{itemize}
    \item \textbf{Normalizzazione Numerica:} Tutti gli identificativi gerarchici originariamente espressi in numeri romani o in forma testuale (es. ``TITOLO I'', ``CAPO IV'') devono essere categoricamente convertiti in numeri interi arabi (1, 4, ecc.).
    \item \textbf{Gestione Dinamica dei Commi:} Poiché nei testi di legge i commi non sono sempre esplicitamente numerati, il modello è istruito a dedurli analizzando le interruzioni logiche dei paragrafi, assegnando loro una numerazione progressiva.
    \item \textbf{Ottimizzazione dello Schema:} Per mantenere leggeri i file JSON, il modello è programmato per non inizializzare chiavi con valori vuoti o nulli; se un livello gerarchico (come un Capo o un Titolo) non trova riscontro nel testo originale, la proprietà viene semplicemente omessa.
\end{itemize}

L'esecuzione di questo processo restituisce una directory di documenti JSON puliti e strutturati gerarchicamente, pronti per essere archiviati e per affrontare il successivo step: l'estrazione delle relazioni semantiche tra le entità legali.

\subsection{Estrazione delle Relazioni Semantiche}

La sola scomposizione del testo normativo in una struttura gerarchica ad albero (Titoli, Capi, Articoli, Commi) rappresenta una condizione necessaria ma non sufficiente per la costruzione di un vero e proprio \textit{Knowledge Graph}. Per sfruttare appieno la potenza topologica di un database a grafo come Neo4j, è indispensabile individuare e mappare esplicitamente i collegamenti e le dipendenze logiche che intercorrono tra le diverse disposizioni giuridiche all'interno dello stesso atto. A questo scopo, la pipeline prevede un secondo step logico interamente dedicato all'estrazione delle relazioni semantiche.

Dal punto di vista implementativo, questo processo si avvale nuovamente dello \textit{Strategy Pattern}. Come orchestrato dallo script principale \texttt{main.py}, una volta conclusa la generazione dei file JSON strutturati, il componente \texttt{FileProcessor} adatta il suo comportamento tramite il metodo \texttt{set\_strategy()}, passando dalla strategia di scomposizione gerarchica a quella di analisi delle relazioni (\texttt{RelationshipStrategy}).

In questa fase, l'input non è più il testo grezzo, ma i file JSON prodotti nello step precedente. Questi file agiscono come una base di conoscenza consolidata e pulita. Ciascun documento viene quindi analizzato dal LLM \texttt{Gemini 2.5-flash}, guidato da un \textit{System Prompt} specializzato (\texttt{gemini-prompt-relationships.md}). Tale prompt è stato progettato per riconoscere i costrutti linguistici del gergo giuridico che indicano un collegamento formale tra le parti del testo. Le relazioni estratte sono:
\begin{itemize}
    \item \textbf{CITA:} Un articolo o comma fa riferimento a un'altra disposizione.
    \item \textbf{ABROGA:} Una disposizione annulla la validità giuridica di un'altra.
    \item \textbf{DEROGA:} Un elemento introduce un'eccezione a una regola generale.
    \item \textbf{SOSTITUISCE:} Una disposizione prende il posto di un'altra.
    \item \textbf{INTERPRETA:} La comprensione di una norma dipende dalla lettura di un'altra.
\end{itemize}

L'output di questa analisi è una rete di collegamenti formali, dove per ciascuna relazione vengono specificati il nodo di partenza, il nodo di arrivo e il tipo di legame. Queste informazioni sono esportate in un nuovo set di file JSON.

Questo approccio a due stadi—prima l'estrazione della struttura e poi l'analisi delle relazioni—permette di separare le responsabilità e ridurre il carico cognitivo del modello a ogni ciclo. In questo modo, si minimizza il rischio di errori e si produce un dataset pulito e affidabile, pronto per essere caricato nel grafo.